% s-258-p258-m001; original PDF 258, 259, 260, 261, 262, 263, 264, 265, 266, 267, 268, 269, 270, 271, 272, 273, 274, 275, 276, 277, 278, 279, 280, 281, 282, 283, 284, 285, 286, 287, 288, 289
\Aphorism{s-258-p258-m001}{352}
In jedem Tropfen schläft ein Weltbild, und das grosse Kunstwerk fällt als Staubbach zu Thal!
\par

Die Seele hat ihre Jahreszeiten; darum muss man sich an seinem Geiste Sommers zu kühlen und Winters zu wärmen wissen.
\par

In jedem Augenblicke treibt nicht das gegenwärtig Reale, sondern das gegenwärtig Imaginäre und Entfernt teste zu künstlerischer Production. Dies gilt für Individuen wie für Völker.
\par

Auf jedes Fremdwort verzichten heisst auch auf jeden Fremdbegriff, jede Fremdempfindung verzichten. Das wollen Einige, — Andere nicht.
\par

Sollte nicht Idiosynkrasie von einem wesentlicheren Ich in uns ausgehen, als unsere bewussten Zu- und Abneigungen? Unsere höchsten Interessen haben alle Ursache, sich nicht der normalen Intelligenz anzuvertrauen; auf Bewusstseinsschienen entgleiste mancher Zug.
\par

Wort ist Werk und Wirkendes, Gedanke Ungestalt und ungestaltend.
\par

Wenn Ihr Geringen nach den Rechten Hoher greift, dann vielleicht — halten sie sie fest! Einem, der abdanken will, muss man nicht den Nachfolger zeigen.
\par

Heute Galgenstrick, morgen Glockenstrang.
\par

Gewohnheit ist schon Verwöhnung; auch das harte Lager muss man wechseln.
\par

Wer für die Vergänglichkeit schaffen will, muss zufrieden sein, wenn es ihm gelingt.
\par

Der gute Stil sucht selbst das nothwendig Seltene und Ungewohnte zu mildern, der schlechte selbst mit dem nothwendig Alltäglichen zu verblüffen.
\par

Der nächste Erfolg jedes Grossen ist Ungerechtigkeit gegen anderes Grosse, die wiedergutzumachen fast einer Rache an jenem Ersten gleichkommt: das Einordnen eines Neuen in unseren Kreis geht nie ohne zwei Gewaltsamkeiten vor sich.
\par

Lasse Kind, Weib und Volk nie wissen, wie sehr du sie liebst!
\par

Einen Gedanken weit zu werfen, bindet die Geschichte den Denker daran; lustig genug, wenn dann der Stein ankommt, das Papier aber unterwegs fortflog.
\par

Ein umfänglicher Mensch sieht seine Peripherie nicht mehr, also gerade dasjenige, womit er Anderen zunächst sichtbar und fühlbar wird. Er müsste, um sich für den Verkehr mit Menschen tauglich zu erhalten, gelegentlich seine Grenzen abgehen, wo bei der grossen Entfernung vom Centrum gewiss Manches der Abhülfe bedarf.
\par

Wo die gröbere Betrachtung noch nicht einmal \Sperr{Zusammenhang} sieht, hat die genauere schon \Sperr{Identität} erkannt, die nicht mehr des Aussprechens werth ist.
\par

Das „Recht zu strafen“ ist ein Unsinn oder ein Pleonasmus: man wähle.
\par

Gesetzt, man dürfe sich der Gesammtheit „Pöbel“ opfern, so doch nie dem Einzelnen innerhalb des Pöbels; dies wäre ja ein Unrecht gegen den übrigen Pöbel!
\par

Der entwickelte Periodenbau der griechischen und römischen Rede ist in unserem eilfertigen Zeitalter, dem die Tugend des Ausredenlassens abhanden kam, nicht mehr möglich. Kein Mensch wagt sich heute noch mit einer ruhigen Syntax hervor, aus Furcht unterbrochen zu werden.
\par

Jemand behauptete, unser Weltall sei aus den Excrementen einer ungeheuren Schildkröte entstanden. Diejenigen, die widersprachen oder sich nicht daran kehrten, nannte er Achelonisten, „Schildkrötenleugner“.
\par

Ein Weib, das geradezu auf die Sinne wirkt, hat den Umweg „Geist“ nicht nöthig.
\par

In jeder leidenschaftlichen Liebe suche ein Gran Perversität; ein normales Gattungswesen \Sperr{kann} sich nicht bis zur Monomanie in den Einzelfall vernarren.
\par

Eine alte Komödie, dass Einer, der dichten wollte, am Bleistiftspitzen scheiterte.
\par

Freundschaft: eine beständige Störung des inneren Geschäfts, ohne die Möglichkeit, dem Störenfried gram zu werden. Darum bedürfen die Schaffenden der grossen Feindschaft; der Freund macht sie unproductiv und entwaffnet zugleich ihren Widerstand gegen das Unproductivwerden.
\par

An einer Unglücksstätte: nicht das Zerstörte, sondern das Erhaltene zu sehen stimmt disharmonisch.
\par

Mitstrebende findet man selten, Mitfaullenzende immer. Will ein Mann sich und seiner Aufgabe entschlüpfen, gewiss hilft ihm — ein Weib.
\par

Der Energie, die ein unreifes Begehren in uns entzündet, werden wir uns bewusst, um zu ernsterem Kampfe auf sie zählen zu dürfen. Unseren Besitzstand zu beleuchten ist auch eine trübe Lampe gut genug.
\par

Die Wegweiser hören auf. Das kann bedeuten: du bist am Ziel — oder du hast es ganz und gar verfehlt.
\par

Der \Foreign{gourmet} des Schmerzes vergiebt sich nie, ein Bitteres, das lebensgefährlich schien, zu schnell verschluckt zu haben. Dieser suchte auf, was er sich nicht ersparen konnte: das Leben wird ihm einmal gute Manieren nachsagen.
\par

Aufs Titelblatt jeder Theorie: Übersetzung in Praxis wird strafrechtlich verfolgt.
\par

Worte sind Manns und Weibs genug, mit einander Gedanken zu haben.
\par

Etwas von dem tiefen Cynismus, mit dem Baubo die Demeter erheiterte, muss einem Weibe eigen sein, das in grossem Stile fasciniren soll.
\par

Was kostet diese That? — Den Thäter. — Das wäre billig.
\par

Im Parlamente Leib ist Gehirn die Opposition.
\par

An Gesagtes wende nicht deinen ganzen Ernst.
\par

Auch im Bette Unglück ruhte sich wohl, liess’ uns nur Hoffnung einschlafen!
\par

Mehr Schnitter her als Wissenschaft, um Welt einzuheimsen!
\par

Lieber befeindet als beweibt; besser, du leidest als deine Sendung.
\par

Das Weibliche am Weibe giebt zu fürchten, sein Männliches zu lachen. Und welch Drittes lieben wir?
\par

Als der Mensch legitim sein wollte, versuchten Natur und Gott nachträglich zu heirathen; Beide waren aber zu alt geworden.
\par

Kürzeste Linie: so compromittirt sich Weltgeschichte nie!
\par

Sich entwickeln heisst das \Foreign{tertium} finden, das auf einer früheren Stufe „non datur“.
\par

Ewig fehlschiessen stimme endlich heiter.
\par

Das Genie wirbelt den Staub nicht auf, sondern macht ihn, wie starker Sonnenschein, nur mehr sichtbar.
\par

Wie wurdest du mich los? — Ich dichtete dich. Jeder Lehrer, das Leben eingerichtet, soll eine Prüfung ablegen, ob er sich darauf versteht, den Schüler zur rechten Zeit fortzuschicken.
\par

Nicht die Geissel ruhn lassen, man habe denn Scorpionen bereit. Moral ist Nothbehelf, für Härteres.
\par

Dinge und Menschen halten einander durch Reibung fest.
\par

Nicht Antworten haben wir über den Menschen hinaus, aber über den Menschen hinaus Fragen: wenig genug, und doch genug.
\par

Die reiche, vornehme Natur erträgt es, Anderen verpflichtet zu sein; das Bedürfniss, Schlag auf Schlag „Revanche“ zu üben, im Guten wie im Schlimmen, ist gründlich plebejisch.
\par

Ein Schiff, das seinem Verfolger entgehen will, fährt auf die Sonne zu: möchten auch wir uns nie anders verbergen, wir Philosophen, als im \Sperr{Übermass} des Lichts!
\par

Das Meer wirft noch Wellen vom gestrigen Südwind, während heute schon Ostwind weht. — Was ist das? ein Gleichniss aus der Leidensgeschichte jedes Genies!
\par

Aus der Befugniss des Mathematikers, reelle Formeln auf imaginärem Wege abzuleiten, macht der Metaphysiker den Unfug, empirische Dinge magisch zu erklären.
\par

Das ist das schlimme Ende des Unrechtthuns, dass man aus allen Näpfen Verzeihung trinken muss.
\par

Eitelkeit: das Diaphragma bei Selbstbeobachtung.
\par

Der Organismus wird vom Centrum zur Peripherie immer weiser, älter, eingespielter. Der Muskel weiss mehr, als der Nerv, der ihn regt.
\par

Theorieen sollen für sich selbst sprechen; ob irgend wer, zum \Sperr{Beispiel} der Verfechter dieser Theorieen, auch nach ihnen lebt und handelt, kommt nicht in Frage. Sieh da! er strauchelt? und doch liegt hier kein Stein des Anstosses! — Nein, er stolperte auch nicht über Etwas, sondern über Nichts. Sein Fuss erwartete die nächste Stufe; konnte er ahnen, dass eure Treppe schon zu Ende sei?
\par

Überhebe dich nicht über deine Mitmenschen, denn der Tod macht alle Menschen gleich! — Dieses Meisterstück von denn erinnert an den Scherz für angehende Mathematiker: einmal Null ist dasselbe wie zehnmal Null (nämlich Null), folglich ist eins gleich zehn! — O diese Tugendlehrer! sie verrechnen sich, sie verschliessen sich, sie verurtheilen sich!
\par

Nur Sonderling sein ist zu wenig; man kann sich von der Alltäglichkeit in horizontaler oder verticaler Richtung entfernen.
\par

Auch das \Foreign{perpetuum immobile} ist noch nicht erfunden, ihr Herren Reactionäre!
\par

Es giebt Argumente, mit denen man keine Zustimmung erzwingen, sondern dem Gegner verrathen will, welche Voraussetzungen man bei jeder Argumentation unangetastet lässt.
\par

Unsere Feinde an sich würden uns wenig anhaben, wenn wir sie nur selbst nicht panzerten, ehe wir sie bekämpfen. Wir denken Gegner und Einwürfe zu Ende, erheben sie ins Ideal und machen uns damit das Leben schwer.
\par

Dass das Alter eines Menschen eindeutig bestimmt ist, kommt nur als Ausnahme vor; für gewöhnlich haben wir verschiedene, ungleichartige, ungleichaltrige Seelen in uns.
\par

Wer die Fabel von der Glückseligkeit des Kindesalters aufgebracht hat, vergass dreierlei: die Religion, die Erziehung, die Vorformen der Geschlechtlichkeit. Wahrheiten sind durchaus nicht immer international. Der Satz „l’essentiel du monde est la volonté“ ist nur im Deutschen eine Wahrheit, dort allerdings eine gestabte, also rechtsgültige.
\par

Allen Menschen darf man zu viel oder zu wenig sein, sich selbst nur — genug.
\par

Wir sehen und denken entzwei, was ursprünglich ganz war; die Dinge haben einander viel lieber als unsere Vorstellungen der Dinge.
\par

Was fragt der Liebende nach Ehe und dauernder Gemeinschaft? Er liebt \Sperr{einmal}, wenn die Gestirnung günstig ist, dann tödtet er, oder stirbt, oder vergisst.
\par

Wer sich grundsätzlich damit befasst, Heuchelei und Schauspielerei an sich und Andern aufzudecken, wird unversehens selbst zum Schauspieler.
\par

In seinen besten Stunden hat der lebende Mensch die Moralität des sterbenden Löwen: er will \Sperr{allein} sein!
\par

Bewusstsein ist nur eine Art, mit den Dingen in Rapport zu stehen, und nicht die intimste.
\par

Ist Reue und Selbstanklage eine vornehme Empfindung? Nur der Plebejer will nach erlittener Unbill des Schimpfens nicht entrathen. Vergebet euren Feinden, auch euch selbst!
\par

Wer Person ist, dem laufen alle Dinge zu. Sie wollen durch ihn zu Worte kommen.
\par

Seine Affecte in Erkenntniss, sein Ganglienleben in Gehirnleben übertragen können heisst philosophiren. Dazu muss man offenbar beide Sprachen verstehen; der rein wissenschaftliche oder der rein pathologische Mensch ist nur ein halber Philosoph, d. h. gar keiner.
\par

Wer nie krank war, weiss nicht, aus wieviel Überflüssigem das Tagewerk des Gesunden besteht.
\par

Es gehört zu den Befugnissen des Dramatikers, dass er, um die Menschen einfach zu halten, die Complication in die Ereignisse verlegt; aber auch nur darum!
\par

Die Bühne der Spiegel des Lebens — aber kein Planspiegel.
\par

Gute Musik und tiefe Glocken sind einer Kirche nützlicher als das klügste Dogma.
\par

Dichter ist, der es allem zu verleiden weiss, was um ihn dichtet: der lauteste Sprecher, der das letzte Wort behält; der übereifrige Anwalt, der dem Vertheidigten in die Rede fällt. Dichter und Leben missgönnen einander, zu Gehör zu kommen.
\par

Man kann den Nexus zwischen zwei Phänomenen längs der Linie causaler Stetigkeit noch so genau kennen: die unwillkürliche Empfindung zieht doch immer die kürzere, die Luftlinie zwischen beiden und glaubt an ein Wunder, an \Foreign{actio in distans.}
\par

Egoist: Einer, der es mit sich gut meint; Altruist: Einer, der will, dass Andere es mit ihm gut meinen.
\par

„Ideale“ sind Speisen, die der menschliche Organismus auf Kosten der Zukunft geniesst; von dieser Seite her ernährt ihn der Wille, wie von der anderen das Wissen. Willenlose Erkenntniss — das heisst, sich durch den After ernähren lassen.
\par

In Dingen des Geistes giebt es kaum einen anderen Reichthum als den an Enttäuschungen.
\par

Das ist ein Original, er meidet die Heerstrassen — aber er tritt auch den einsamsten Waldpfad zur Heerstrasse breit. Vielleicht hat er für ein Original zu plumpe Füsse.
\par

Mancher lässt sich um seinen Schmerz betrügen und glaubt damit noch das Glück erfunden zu haben.
\par

Es ist billig, dass wir unsere Erfahrungen theurer bezahlen als Andere; für uns, die wir zwischen den Zeilen lesen, steht mehr darin! Geistige Männer sind oft für einander stumpf, unempfindlich, unempfänglich — eine traurige Geschichte. Eine Flamme wirft Schatten im Lichte der anderen, als wäre sie ein schwarzer Körper, und doch sind sie Schwestern!
\par

Wer wird aus den Trümmern eines grossen Glücks die Überbleibsel zusammenlesen! Wer ganz reich gewesen ist und es wieder einmal sein wird, ist inzwischen stolz genug zur ganzen Armuth!
\par

Hat man einmal sein Zartgefühl zerrissen, so schaffe man sich etwas Besseres an: es wächst nicht wieder zu.
\par

Wenn Schweigsame einmal zu reden kommen, so schütten sie mit einem Augenblicke zehn Jahre vor- und nachher aus. Darum nehmen sie es so tödtlich übel, nicht verstanden zu werden: sie haben zuviel auf eine Karte gesetzt, sich für zulange ausgegeben.
\par

Wer die Rache schilt, vergisst, dass es Reinheiten giebt, sich von ihr zu waschen.
\par

Lenkungen sieht man’s zu spät an, wovon sie ablenkten, zu früh, welchem Nichts sie zulenken.
\par

Das Glück ist schnell bereit, Verschuldungen wieder gut zu machen: dabei aber stösst es auf den ganzen Trotz, dessen ein Zurückgesetzter fähig ist, und wird selber trotzig.
\par

Was liegt an Menschen und Dingen! Es wird immer \Sperr{grosse} \Sperr{Augenblicke} geben.
\par

Kniee vor einem Stein: sofort stellt sich ein Gott darauf. Stumpfe deinen Blick: sofort siehst du ein Räthsel. Greife müde: sofort greifst du auf Widerstand. Säume zu gebieten: sofort verbietet dir Sitte.
\par

Schon eine Raumdimension mehr erklärt Wunder über Wunder: um wieviel mehr erst eine qualitative Weltsteigerung! Nicht einmal die Philosophie des Plebejischen ist zu bekämpfen, geschweige denn der Plebejer selbst.
\par

In dem Augenblicke, wo uns eine Täuschung gelingt, reizt es uns schon wieder sie zu zerstören.
\par

Ursache muss festes Gewissen haben, um alle ihre Wirkungen zu verantworten.
\par

Höchster Grad von Freundschaft: einem Überflüssigen zuliebe, damit er sich aushalte, selber überflüssig werden.
\par

Begründet, erfindet, erlügt Unterschiede — aber verwischt keinen! Besser eine utopische Rangordnung als die plausibelste Gleichheit.
\par

Hör’ auf, mich zu lieben, dass du wieder gut zu mir sein kannst!
\par

Nach tiefer Erschütterung lange schweigen, daraus wächst die grosse Kunst. Man wird an den Dingen nie dann zum Dichter, wenn man zu früh über sie Worte findet.
\par

Das Heiligthum eines Aufrichtigen: eine Galerie von Menschen und Landschaften, denen er \Sperr{nicht} begegnet ist.
\par

Wer noch hofft, hat die Hand nicht frei zum Experiment.
\par

Verdienst du die Strenge, deren du bedarfst? Nicht eher, als bis du selbst sie dir erzeigen kannst.
\par

Schone Alle, von denen du einmal nicht geschont sein willst.
\par

Weiss er die Stange nicht als Knüttel zu brauchen, ist der Fahnenträger verloren.
\par

Auch die Pausen sind Musik, auch Nichtgesagtes gilt es zu hören.
\par

Fehlgeschossen! — Ja, Freund, warum zieltest du auch dahin, wohin du treffen wolltest?
\par

Hat man sich über kleines Missgeschick geärgert, lerne man über die grosse Irrationalität lachen, dass Jeder, der plant, nur Eines vergisst: die Hauptsache. Umwege von Viertelstunden mögen verdriessen; aber mit dem ersten Schritt nach Westen gerathen, wenn man nach Osten will, — dabei giebt es nichts zu grollen.
\par

Mit der Umarmung hört es auf, dass man einander die Hand geben kann.
\par

Was zu Zweien möglich ist, muss man wissen, um Einer zu werden.
\par

Liebe und Sinnlichkeit sind zu verwandt, um Freunde sein zu können.
\par

Ein Siebenmeilenschritt vorwärts, jeder folgende zurück: Normalverlauf des altruistischen Anfalls, genannt Liebe.
\par

Lass kleine Weiblein über zerstörte Illusionen klagen: du kamst nicht auf den Markt, um wohlfeil einzukaufen.
\par

Nahm ich für dich Partei: weh dir, fällst du nun von dir ab!
\par

Ad absurdum geführte Vernunft ist lammfromm gegen lügengestrafte Nerven.
\par

Gott wird dauern: Atheismus conservirt ihn.
\par

Am Erlebten rächt uns Darstellung, oft über gerechtes Mass.
\par

Alles vergeben: eine sublime Art, zu Bosheit zu reizen.
\par

Nicht in Etwas sein Gegentheil, sondern in Etwas Nichts entdecken ist die verletzendste Enttäuschung.
\par

Die \Sperr{Seele} vermag es noch nicht, Selbstberührung von Fremdberührung auch im Dunklen zu unterscheiden: sie fühlt Concubinat und ist Morgens abergläubisch erschrocken, allein zu liegen.
\par

Gott ist ein Bedürfniss nach äusserster Ungenirtheit; Menschen mit Disciplin haben dergleichen nicht nöthig.
\par

Der freie Geist will nicht Herrenlosigkeit, sondern centralisirte Botmässigkeit; er verlangt, dass die niederen Imperative, die seine Lenkung in die Hand nehmen, selbst wieder höheren gehorchen.
\par

Als letzte Versuchung überwinden wir den Hochmuth, uns versucht zu glauben.
\par

Der Geräderte hat wenigstens Musse, still zu liegen — vielleicht das erste Mal in seinem Leben.
\par

Beim ersten Besuch legt man gewisse Garderobe-Stücke \Sperr{nicht} ab. Es ersparte manches Peinliche, wenn man sich z. B. vornähme, bei der ersten Liebe seinen Geist anzubehalten.
\par

Das beste Gespräch führen Zwei über Zweierlei; unmöglich, \Sperr{nicht} am Andern vorbeizureden.
\par

Schriftsteller, die von ihrer Kürze sprechen, lassen sie gewöhnlich vermissen. Die wirklichen Brachylogen halten sich immer noch für zu redselig.
\par

Irgendwem zu begegnen, gehe man in sein Verbotenes.
\par

Es erfüllt mit Vertrauen zur Menschheit, dass sie es nicht erträgt, ohne régime zu leben; es ist ihr \Foreign{horror vacui.}
\par

Unter den Bewunderern entgleister Grösse ist auch der Philister; hier, wo es keine Gefahr mehr bringt, entschliesst er sich gern auch Geist zu haben.
\par

Das allgemein Geglaubte giebt eine krumme Richtschnur, mag man sich nun danach oder dagegen richten.
\par

Iss auf, was du hast: dann \Sperr{kannst} du verhungern, \Sperr{kannst} folglich etwas Neues zu essen bekommen. Man soll, so oft es geht, seine Existenz von Dingen abhängig machen, deren Existenz noch unbewiesen ist: das sicherste Mittel, ihr Entdecker zu werden. Als „verlorener Pfeil“ ist das Leben immer noch wohl angewandt.
\par

Jeder Mensch hat der Menschheit mehr zu verzeihen als sie ihm.
\par

Gegen sein Wissen glauben stärkt die Persönlichkeit, wie gegen den Strom schwimmen die Arm- und Beinmuskeln; aber derlei Sport will nicht übertrieben sein.
\par

Der Philosoph hofft, wenn er sich ein paar naturwissenschaftliche Facta aneignet, man werde seine ganze Philosophie für Wissenschaft halten: er gleicht dem Panoramenmaler, der einige körperliche Gegenstände in den Vordergrund stellt, um auch für seine bemalte Leinwand den Schein der Körperlichkeit zu erwecken.
\par

Die unmittelbare Gegenwart ist trivial, erst durch Zukunft oder Vergangenheit nehmen die Erlebnisse Grösse und Grossartigkeit an.
\par

Man darf Analyse und Allgemeinheit nur bis zur vorletzten Stufe treiben, wenn man die Dinge nicht entfärben und entkörpern will.
\par

Wenn der Mensch producirt, hat er die Fühlung mit der Aussenwelt verloren; statt seiner repräsentirt ein Ich zweiter Klasse, ein Ego adlatus.
\par

Im Theater des Lebens giebt es viele wirklich sehr schlechte Plätze. Aber das aufgeführte Stück ist gut.
\par

Grosser Schmerz ist nur geliehen: hast du ihn ausgelebt, so gieb ihn zurück.
\par

Es lohnt sich nicht, auf seine Tageseintheilung Geist zu verwenden, wenn die Woche weniger als sieben Sonntage hat.
\par

Der Philister hat alles gute Gewissen aufgekauft, wahrscheinlich um den Preis dafür in die Höhe zu treiben. Und wirklich, mancher höhere Mensch gab sein Letztes (sich selbst!), um sich wieder einmal moralisch zu fühlen.
\par

Ein vornehmer Geschmack schützt auch vor vielen Arten der Unmoralität, aber nicht desshalb ist er vornehm.
\par

Für ein fertiges Publicum schaffen, in eine fertige Moralität hineinhandeln ist weder Kunst noch Sittlichkeit ersten Ranges. Was der Mensch in die Schauläden seines Bewusstseins legt, ist fast immer Trödel, aus dem man nicht auf den „innern“ Menschen schliessen darf.
\par

Wer sich nicht halten kann, soll nicht warten, bis er fällt, sondern hinabspringen, „sich opfern“. Es sieht so besser aus! mehr will ich nicht sagen.
\par

Eine Wahrheit, die das Volk „begriffen“ hat, trägt die Spuren schmutziger Finger.
\par

De omnibus dubitandum: also auch an der Verbindlichkeit dieser Vorschrift!
\par

Erkenntnisse, die wohl oder wehe thun, sind mit unwissenschaftlichen Zusätzen vermischt; die reine Wahrheit redet keine Gefühlssprache.
\par

Wer nicht \Sperr{sich} will, soll nicht schlechthin Anderes wollen, sondern Besseres.
\par

„Und“ ist das \Greek{πρῶτον} \Greek{ψεῦδος} der Sprache.
\par

Sechs Tage \Sperr{sollst} du arbeiten, der siebente \Sperr{ist} ein Ruhetag und Gott geweiht! So ist es recht; otium ist eine Kunst und lässt sich nicht gebieten.
\par

Optimist darf man nicht auf einen Theil des Daseins, sondern nur auf das Ganze hin sein.
\par

Lieber ans Kreuz genagelt sein, als in einer peinlichen Lage sich noch selber festhalten müssen.
\par

Die moderne Bildung: Acteure und Redacteure.
\par

Er schweigt viel, aber glaubt ihm nicht! er verschweigt nichts! er würde verstanden werden, wenn er redete!
\par

Die wohlfeilen Arten der Erkenntniss genügen Dem nicht, der für Erkenntniss mit seinem Glück und Leben zahlen will.
\par

Der Krieg antwortet auf mehr Fragen, als gestellt wurden.
\par

Um alles, was Object heisst, wider dich zu haben, werde Solipsist! Dem es beizubringen, dass sie auch existiren, ist aller Objecte Lust.
\par

Wer gerichtet hat, richte nun sich — wenn aus keinem andern, dann aus dem Grunde, weil er das Beil einmal in der Hand hat.
\par

Man muss auch sich selbst die nächsten Wünsche an den Augen absehen. Einer, der gut mit sich auskam, sagte: Ich bitte mich um Nichts, was ich mir nicht längst gewährt hätte.
\par

„Eine Scene machen“ drückt vortrefflich aus, dass man dazu theatralisch angelegt sein muss; grosse Seelen zürnen still!
\par

Das Haupt, das die schwerste Krone trägt, kann nicht mehr nicken, nur noch langsam verneinen \Dots{3}
\par

Weil du nicht Pygmalion sein kannst, muss dich der Hammer am Meissel rächen?
\par

Mitleid ist ein Quiproquo, das Dem schmeichelt, der dabei gewinnt.
\par

Ideale ohne Idealisten würden viel mehr Zuspruch haben.
\par

Dass sich die Menschen an Dingen, die sie nicht verstehn, die Köpfe zerbrechen, ist nicht das Ärgste; wenn sie nur nicht mit diesen Köpfen weiter dächten!
\par

Es gehört mehr Mensch dazu, Errungenes zu geniessen, als es zu erringen.
\par

Die ausschweifendsten Vorstellungen vom „Recht der Jugend“ hat man zu der Zeit, da man am wenigsten jung ist: in den zwanziger Jahren.
\par

Mit seiner Vollkommenheit thut man Anderen nothwendig wehe: und die Meisten suchen sie gerade darum.
\par

Eine Gegenwart, die rasch und heimlich genossen werden muss, hinterlässt den feurigsten Nachgeschmack.
\par

Es liegt höhere Menschlichkeit darin, dass man so Vieles auf Erden für Geld haben kann. Wer die conventionellen Formen des Erstehens und Sichabfindens ersann, wusste um eine vornehme Seele und deren Bedürfnisse; nur der Pöbel liebt, hasst, bittet und dankt mit der ganzen Person.
\par

Geist und Tugend vertragen sich schon in einer Seele, nur müssen sie von Tisch und Bett geschieden sein.
\par

Wiederholung ist die mildeste Form von Parodie.
\par

Wir sprechen immer Höheres aus, um darüber hinaus noch Höheres zu verschweigen; so wachsen wir mit unserem Sagbaren und Unsagbaren in die Höhe.
\par

Menschliche Eitelkeit möchte sich überreden, der Sprung vom Thiere zum Menschen sei so ungeheuer, dass ihn die Natur nur einmal gemacht haben könne. Aber die hat lange Beine!
\par

Wahrscheinlich sind am Denken Innervationsimpulse des Kehlkopfs betheiligt. Wir denken nicht nur in Worten, sondern geradezu in Klängen, in rhetorischen Farben und Stimmungen — wir stellen uns auf eine fingirte Rednerbühne vor fingirte Zuhörer.
\par

Gespensterfurcht ist ihre eigene Enkelin: Mutter und Tochter der Religionen.
\par

Noch ein Jahrhundert den Körper studirt, und Niemand liest mehr die Übersetzung „Geist“.
\par

Eingesperrt in der eigenen Haut: ist’s ein Wunder, dass der Gefangene närrisch wird, dass der Mensch nach dem Kerkerwechsel „Liebe“ schmachtet?
\par

Weltuntergang wäre heute noch gute Ernte für Wirthe und Kuppler: vertagen wir ihn, bis er Besseren Zulauf schafft.
\par

Man erzählt mit falschen Accenten, solange man von jeder Gattung nur Ein Erlebniss hinter sich hat.
\par

Gegen die Liebe predigt Ohnmacht mit ihr umzugehen. Eine Naturkraft wünscht man nicht hinweg: man zwingt sie für uns zu arbeiten.
\par

Danke deinen Irrthümern, wenn du so alt wirst, Wahrheiten zu erleben.
\par

Hart gegen Andere ist die Vorschule zu Hart gegen sich, also \Foreign{prae limine} zu billigen.
\par

Wer sich berufen fühlt, Illusionen zu zerstören, handelt selbst aus Illusion.
\par

Zeus hat von den Griechen eine bessere Meinung als sie von ihm; die Griechen erfinden Götter, um gut von sich zu denken. Das Christenthum erst bringt dies Verhältniss ins Reine — wirklich ins \Sperr{Reine}?
\par

Den Kreuzestod sterben für seine Wahrheiten — gut! Aber ein Anderes ist, sich um ihretwillen von Würmern und Giftfliegen benagen lassen. Moderne Martyrien sind unästhetisch, darum geht man ihnen mit Recht aus dem Wege.
\par

Die letzten Ausläufer einer Gedankenkette sind entweder zu spitz oder zu niedrig, um die Schneekrone der Wahrheit zu tragen.
\par

Liebe ist oft eine Verlegenheit und Ausflucht Derer, die sich nicht auf die kälteren Umgangsformen verstehen.
\par

Alles Schaffen geschieht im Rausche und wird weder von den nüchternen Menschen noch von den nüchternen Augenblicken des schaffenden Menschen verstanden.
\par

Der starke Glaube erübrigt die Existenz des Geglaubten.
\par

Der wirkliche Mensch sieht den idealen über sich schweben und begreift nicht, dass es die Schwere der \Sperr{Realität} ist, die seine Wagschale zum Sinken und die des Anderen zum Steigen gebracht hat.
\par

Religion ist zu bequem für uns, aber der heutige Mensch ist zu bequem zur Religion. Es giebt Freunde als Nothhelfer und Freunde als Genusshelfer, Mitleidende und Mitfeiernde; die zweite, edlere Art bleibt einander fern, solange sie am Alltag des Lebens zu schleppen hat, und schaut nur von den selten erreichten Gipfeln der Reife, Ruhe, Meisterschaft zu ihresgleichen hinüber.
\par

Sich mit zahlreichen Forderungen von Sitte und Sittlichkeit umstellen ist das beste Mittel gegen Langeweile; der Tugend bringt der ganze Tag Besuch — von „Versuchungen“.
\par

Durch vieles Denken wird man gegen alles Denken misstrauisch.
\par

Die meisten „Interessen der Menschheit“ sind verkappte Diätfragen.
\par

Seid Welt, und lasst Andere Welt erklären!
\par

Das Ross That ist gestürzt: wie sie herangekrächzt kommen, die Aasgeier Reue und Ruhm!
\par

Du zweifelst? Kein Glaube war mehr Glaube, als dein Zweifel.
\par

Hülle verkehrt mit Hülle, die Fülle bleibt stumm, also unbeweisbar.
\par

Der Gekettete erkennt, wie oft er sich ehedem mit Fäden binden liess.
\par

Der Geduld eines Engels kommt nur eine wirkliche nahe: die des Raubthiers, das auf Beute lauert.
\par

Ärmster der Verkünder! Alle hast du überzeugt — ausser dir!
\par

Befehdung eines Ausser-mir ist Befreiung von einem In-mir.
\par

Dort sinkt mein Glück\Dots{2} Sah ich es je so schön?
\par

Draussen hängen die Dinge fester an einander als in unserem Kopfe.
\par

Allen erlaubt, mir verboten: so redet der freieste Geist. Misslingen im Kleinen ist noch lange keine Berufung zu Grossem!
\par

Ein Gedanke sucht Obdach. Ich vermiethe mich ihm — für heute!
\par

Gedanke beim Heirathen: was ich vor mir nicht schützen konnte, muss ich vor Schlimmerem schützen.
\par

Um den gefährlichsten Problemen zu begegnen, muss man nur ohne Jagdgewehr ausgehen.
\par

Jeder Erfolg ist Ablenkung, wäre nur auch jeder Misserfolg Zurechtweisung!
\par

Grösse ist, sich ihrer begeben können.
\par

Die Frau überwintert, der Mann bringt zur Reife.
\par

Der Mensch wird böse, wenn man ihm wehrt, auf seine Art zu Schaden zu kommen.
\par

Auch bei der Selbsterziehung lässt es sich so einrichten, dass der Geführte Führer zu sein glaubt.
\par

Der Astronom bringt sich um den Himmel, der Historiker um die Helden, der Ehemann um die Geliebte. Man trinkt sich nüchtern, handelt sich ärmer, lebt sich todt.
\par

Nimm mein Leben! aber keinen Heller Lösegeld! — Alles opfert sich leichter als Weniges.
\par

Man heizt mit Liebe, um den Geist in Gang zu setzen; gäbe diese Kohle nur nicht so erbärmlich viel Qualm!
\par

Lucrezia Borgia, das Weib ohne Gedächtniss, \Sperr{folglich} das einwandfreie Weib! Man hat immer erst eine Vergangenheit, wenn man sich nach ihr umsieht.
\par

Sind Worte nicht dumme Vögel, dass sie auf die Leimruthe „Gedanken“ gehen?
\par

Es ist schade, dass man die Rechtmässigkeit einer Fragestellung nie früher anzweifelt, als bis tausend Versuche der Beantwortung missglückt sind. Wir könnten sonst in der gleichen Zeit die doppelte Culturstrecke zurücklegen, aber „die Menschheit soll nicht so schnell zum Ziele“.
\par

Was der Philosoph für\Sperr{ unmittelbare Gewissheit} ausgiebt, ist gewöhnlich eine ganz verzwickte Lüge.
\par

\Foreign{Credo quia absurdum} — so redet kein wahrer Glaube. Auf diese Weise schriebe ja doch die Vernunft vor, was geglaubt werden muss, nur \Foreign{per contrarium}?
\par

Der Affe als degenerirter Mensch ist vielleicht verständlicher denn der Mensch als progenerirter Affe.
\par

Ungerechtigkeit ist das Recht der Schaffenden, Widerspruch das der Realitäten.
\par

In Luft und Wasser nimmt der Druck der Umgebung von oben nach unten zu, in der menschlichen Gesellschaft von unten nach oben.
\par

Jeder grosse Mensch hat das, was seinen Nachbetern fehlt, das Gewissen und die schreitende Folgerichtigkeit seiner Entwicklung; desshalb wird er oft von ihnen scheinbar überholt.
\par

Selbständigkeit will nicht gesucht sein. Man soll sie haben, aber man soll eher noch nach Führern und Vorbildern suchen.
\par

Wenn selbst der Pöbel liest, schreibt und discutirt, ist lange kein Krieg zu erwarten. Die Bauern sagen: Bergaufweiden des Viehs gilt als Zeichen guten Wetters.
\par

Wer bisher die Menschen befreien wollte, war selbst noch zu sehr Knecht.
\par

So will’s die Logik! — Wenn weiter Niemand!
\par

Das ist Philosophie, die versagt, wenn man ihrer bedarf; der Nothhelfer giebt es schon zuviele.
\par

Man kann die Geliebte zum Weibe, zum Kinde, zur Mutter haben wollen: auf das Merkmal „Generation“ ist in der Liebe zu achten. Seiner Misserfolge beim Publicum rühmt man sich, warum nicht auch derer beim Weibe?
\par

Nicht alle Seelenkämpfe sind gefährlich; man lasse sich von den Dichtern nichts einreden.
\par

Der berufene Mensch nimmt es mit den Anlässen, sich zu opfern, nicht genau genug; wo eine Thür zum Untergange sich aufthut, meint er, es gelte ihm.
\par

Von Jemandem nicht scheiden können, ohne ihm „die Wahrheit zu sagen“, ist Sache der gelbgrünen, galligen Seelen. Ausserdem schneidet es die Wiederanknüpfung ab, und derlei sollte man sich offen halten, selbst wenn dieser Jemand das \Sperr{Leben} ist.
\par

Wie würde in einigen Generationen die Menschheit aussehen, wenn Sitte geböte, am Tage, bei hellem Sonnenschein, zu zeugen?
\par

Gerade der Mann mit Tiefgang kommt schwer über die Sandbank Weib hinweg.
\par

Der Denker, der liebt, muss sich als zurückgelegtes Stadium, als Larve conserviren; er hat zu schonen — nicht sich, sondern sein Bild in einem vertrauenden Herzen.
\par

Heute schwöre! welcher Fahne, erfährst du morgen.
\par

Es giebt viele Moralen, wohl uns! So schläft der Strom und Windzug Immoralität nie ein, der sie ausgleicht! Frisches Wasser, frische Luft für die Staub-erstickten, die Stubenhocker einer einzigen beschränkten Sittlichkeit!
\par

Ein explicite sexuelles Dichten ist unreif und unerträglich. Das Geschlechtliche mag überall als Grundton hörbar sein, aber nicht für sich allein herausklingen; es muss ins Leben resorbirt werden wie das Sperma ins Blut.
\par

Die schlichte Ausdrucksweise ist stärker, anspruchsvoller als die hyperbolische; vielleicht muss man gerade, um nicht zu übertreiben, in Superlativen sprechen, d. h. zu einem gewissen Grade des Misstrauens auffordern.
\par

Nach mir die Sintfluth! so denkt jede Liebe und verrechnet sich wie jener Verschwender, der sein Vermögen überlebt. Vielleicht ist Liebe ihrer Folgen wegen da: als Anreiz, sich einmal ganz auszugeben und dann eine Weile von Nichts zu leben.
\par

Was ich im Rausche gethan, lasse ich am wenigsten im Stich, — sagte ein Nüchterner.
\par

Der fatalistischere Fatalismus ist, der die Dinge \Sperr{nicht} gehen lassen will, der jeden Windstoss als Signal zum Curswechsel ansieht.
\par

\Sperr{Grenze des Charakters.} Will man auch nur Denen Recht geben, die wirklich Recht haben, so kann man nicht Ein und derselbe Mensch bleiben.
\par

Warum in Feinden wählerischer sein als in Freunden?
\par

Die glücklichste Liebe ist die muthigste, die sich auf halbausgebranntem Vulkan anbaut.
\par

Ein zurückhaltendes, schweigsames Weib: was hält es eigentlich zurück, was verschweigt es denn? Ich lobe mir das kluge, heitere, zutrauliche Weib: das einzige, das aufrichtig ist und sich nicht inscenirt.
\par

An der Grenze zweier Erlebnisse: all mein Erworbenes lasse ich, um Zoll zu sparen, drüben und fange im neuen Lande von Neuem an.
\par

Es sind nicht einmal die eigenen Tugenden, mit denen wir uns das Leben erschweren.
\par

Verdorbenheit rechnen wir einem Weibe hart an, wenn wir es nicht verdorben haben.
\par

Schade! auch auf ewige Antipathie ist kein Verlass!
\par

Erfand ein erhörter oder ein verschmähter Vogel das Singen? Ist die Kunst Dank oder Rache? Die Irrthümer, von denen wir wieder „zurückkommen“, sind damit nicht entwerthet. Bereuen wir Reisen, weil wir einmal nach Hause müssen?
\par

Dass, trotz seiner grundverkehrten Erziehung zur Ehe, das Weib in der Liebe weniger phantastisch ist als der Mann, dankt es seinem unzerstörbaren „Genie der Gesundheit“, um dessen Mitgenuss wir uns in aller Demuth zu bewerben hätten, statt unsererseits das Weib romantisch zu machen.
\par

„Brennende Fragen“, welch ein Widersinn! Wenn sie wirklich brännten, würdet ihr schon die Antwort wissen!
\par

Dem Geliebten verzeihen wir sogar seine Vorzüge.
\par

„Sein Blut komme über uns und unsere Kinder!“ — Aber wer wird Unmündige beim Wort nehmen?
\par

Liebe ist Loskauf von Güte und Dank; Steine statt Brotes.
\par

Worte über ein Erlebniss, einen Menschen finden heisst einen Sarg zunageln.
\par

Das Leben schmeckt mir gut und — schlecht genug, jedenfalls bedarf ich keiner Würze und Nachhülfe und künstlichen Piquanterie, ihr religiösen Muskatköche! Giftmischer! Sinnenkitzler!
\par

Walten ist Gewalt, jedes Thun ein Unrechtthun. Gott kann ein Schöpfer sein, oder ein Gerechter: nicht Beides.
\par

Was dem Auge wohlthut, würde die greifende Hand verletzen: starke Formen, Krümmungen, Spitzen, Kanten. Unser „Gefühl“ will die Dinge glatt, flächenhaft, rund, „ideal“, also dem Willen unseres Gesichts entgegen.
\par

Absolutes Herrschen ist Dienen, wie es ein Zwang ist, sich ungezwungen zu bewegen. Der Mensch will hier gehorchen, dort befehlen; nur gehorchen, nur befehlen ist ihm unnatürlich. Die Grundsätze richten sich zehnmal nach dem Handeln, ehe einmal das Handeln nach Grundsätzen.
\par

Galiani über Rabelais: \Foreign{Il ressemble au cul d’un pauvre homme, frais, dodu, sale et bien portant.}
\par

Je besonderer wir geformt sind, desto seltener finden wir uns congruente Erlebnisse, und desto fester haften sie in uns.
\par

Wer viel Tugend hat, braucht viele; wer sich nicht in Versuchung führt, verträgt keine.
\par

Über Religion, Staat, Ehe denkt man ungefähr so: es ist schon viel, dass diese Dinge nicht an ihrer eigenen Praxis zu Grunde gehen; aber noch dazu theoretische Beleuchtung — die vertragen sie keinesfalls!
\par

Die holde Scham, die nicht ätzt: so schämt sich Wohlgerathenes seiner Reife, das Weib seiner Mutterschaft, die Welt ihrer Schönheit.
\par

Der Weg zum Können heisst Müssen, nicht Dürfen.
\par

Nonum prematur in annum: aber nicht das Gedichtete, ehe es mitgetheilt, sondern das Erlebte, ehe es gedichtet wird.
\par

Erst das Vergessene hat voll gewirkt.
\par

Sich \Sperr{ver}lieben: eine Freimüthigkeit, die der deutschen Sprache Ehre macht.
\par

Verschwinde du! oder dein Werk ist verschwendet!
\par

Man spricht für eine Sache, was man dem Thun für sie entzieht.
\par

Wenn man den Idealen nicht Gelegenheit giebt, sich zu Tode zu practiciren, so wundere man sich nicht, dass sie ewig leben.
\par

Du bist noch Fleisch? So ist deine Lehre noch Wort.
\par

Schäme dich, dass dies Unglück dir aus dem Wege ging! Was mit uns gewollt wird, verräth sich als Widerstand gegen das, was wir wollen.
\par

Was Jungfer Ideal dir geboren, — lass es Frau Wirklichkeit nicht sehen!
\par

Gegenliebe ist ein Todesstoss, den kaum die gesündeste Liebe verwindet.
\par

Wem zuviel gegeben worden, dem wird das Überflüssige wieder abverlangt werden. Und gerade das Überflüssige entbehrt sich am schwersten.
\par

Das Leben ist ein Umweg zwischen Geburt und Tod, der \Sperr{subjectiv} nähere Weg führt vom Tode vorwärts durch das Unendliche.
\par

Polemische Zeiten und Menschen sind arm an gründlichen Gegensätzen.
\par

Manche begehen Haufen Unrecht, um nicht eine Handvoll begangenes einzugestehen.
\par

Zwang ist auch Sicherheit, Gefahr schirmt vor Gefahren.
\par

Jedes Temperament ist eine Fälschung, die dem Leben Sinn oder Unsinn, Schönheit oder Hässlichkeit zulegt, — zulügt.
\par

Die Moral muss ihren eignen Namen vergessen haben, da sie überall, wo von Werth und Schätzung die Rede ist, sich und ihr Urtheil angerufen glaubt.
\par

Man muss als Seele so hoch und weit gebaut sein, dass der Lärm im Erdgeschoss nicht bis in die oberen Stockwerke dringt.
\par

Kurz und gut: was will ich eigentlich vom Leben? — Mich!
\par

Das Leben überwinden ist eine Redensart für den Thatbestand: vom Leben überwunden werden.
\par

Wer sich an der „ewigen Gerechtigkeit“ des Daseins unbedenklich ergetzt, ist weder Richter noch Verbrecher, sondern gaffendes Tribünenpublicum.
\par

Dunkle Ahnungen und Dränge peinigen; daher die Melancholie der Jugend und die Heiterkeit des nicht gefühlsärmeren, aber gefühlsklareren Mannesalters.
\par

Den Einen zerhaut das Leben mit scharfer Schneide, den Anderen zersägt es mit Luntenstricken. Wer betet, sollte sich Marterung mit solchen Instrumenten erbitten, die wirklich und ursprünglich zum Martern bestimmt sind.
\par

Es giebt Fragen, in denen man ohne einigen Geist gar nicht irren kann.
\par

Honig und Stachel der Erlebnisse liegt vor und hinter ihnen, nicht in ihnen.
\par

Es ist Philosophenlust, über Dinge reden zu dürfen, die ihn nichts angehen.
\par

Der Mensch ist zu gut für diese Welt, also wird er in jener weiterleben. — Diese Welt ist zu schön für den Menschen, darum sollte er in dieser weiterleben, noch einmal und hundertmal leben.
\par

Weibesschönheit ist Wille zu gebären, Mannesschönheit Wille zu zeugen: nur mit dem Unterschiede, dass der Mann auch auf andere als leibliche Vaterschaft hinauswill.
\par

Um uns selbst besser zu schmecken, spülen wir uns den Mund aus — mit Gesellschaft.
\par

Man soll Wahrheiten sagen, wie man einen Stein schleudert, nicht wie man etwas aus der Tasche verliert.
\par

Jeder Philosoph hat einen Gedanken, dem gegenüber er aus der Rolle des Denkers fällt.
\par

Dass uns eine Erkenntniss juckt oder kitzelt, beweist, dass wir nicht fest genug zugegriffen haben.
\par

Ehe man seine Moral Anderen auflädt, soll man sie selbst eine Lebensstrecke lang getragen haben. Eine nicht gar zu wissenschaftliche Wiedergeburtslehre würde den Selbstmord gestatten müssen, etwa als Abkürzung eines Besuchs, der den Hausherrn bei schlechter Laune trifft, mit dem Vorbehalt, zu gelegenerer Zeit wiederzukommen.
\par

Niemand hält sich gern für überflüssig, — lieber noch für schädlich.
\par

Wir nennen den Löwen ein Raubthier und die „lieben Vöglein“ unschuldig; Fliegen denken umgekehrt.
\par

Beim Bau einer Weltanschauung vergisst man gewöhnlich die Fremdenzimmer.
\par

Es ist eine Pein für tiefe Menschen, dass die flachen Erklärungen meist die richtigen sind.
\par

Das Object — eine Übereinkunft, ein Vertrag vieler Subjecte.
\par

Motivation ist ein ästhetischer Fehler; die Dinge im Kunstwerk müssen ihr Woher und Warum am Leibe tragen, nicht auf dem Redekarren hinter sich herziehen, — und müssen überdies so berechtigt aussehen, dass man die Leibesvisitation vertrauensvoll unterlässt.
\par

Der Mensch muss einen Herren haben; was folgt daraus? „Ich muss sein.“ Nein, du musst haben!
\par

Das Leben lieben ist Mannesvorrecht; für schöne, reife und etwas cynische Weiber haben Jünglinge kein Verständniss.
\par

Ein Problem in aller Eile zu formuliren, damit es nur als Problem nicht davonlaufe, muss man im Ausdruck nicht wählerisch sein. Gestatte uns einstweilen, dich als Monstrum zu bilden: warum bist du Sphinx!
\par

Der theatralische Cyniker, der sein Licht (seine Diogeneslaterne) leuchten lässt, wie würde er in Verlegenheit gerathen, wenn ihm der Mensch begegnete, den er zu suchen vorgiebt! Durch seinen Gott hat der Mensch den Weg vorwärts gefunden, durch seinen Gott könnte er leicht denselben Weg wieder zurück finden. Wir sind nicht zuverlässig genug, als dass wir nicht nöthig hätten, diese Brücke hinter uns abzubrechen.
\par

Die christlichen Märtyrer starben für die Sache Jesu Christi, und Jesus — desgleichen: um so viel steht er über ihnen, wie es schwerer ist, für die eigene Sache und für eine unbewiesene Zukunft zu sterben als für eine angefochtene Gegenwart, eine blosse Partei.
\par

So viele Menschen sind eine Beleidigung der Menschheit, aber Niemandem fällt es ein, sie zurückzunehmen! Natur hat das \Foreign{point d’honneur} der Raufbolde: provociren und nicht revociren.
\par

Es ist verfeinertes Wehethun, sich das passive und Anderen das active Unrecht aufzubürden.
\par

Der Wunsch, beachtet zu werden, ist der stärkste Trieb der Menschheit; sie will nicht glauben, dass in unserem Erdenwinkel gelebt werden kann, ohne dass das gesammte Weltall Ohren und Augen aufsperre. Wir fühlen uns immer noch als Schauspieler und schneiden die wunderlichsten Gesichter; eines Tages werden wir ein sehr dummes machen, wenn uns zum Bewusstsein kommt, dass wir kein Publicum haben.
\par

Das Leben knüpft und löst Knoten, die das Denken weder lösen noch knüpfen kann: vielleicht ist das Leben um eine Dimension höher!
\par

Auf Nachwelt rechnen sollte durch Vergessen bestraft werden, und Der nur überleben, der willig in Vergessen sank.
\par

Auch das Leben sucht seine Eitelkeit darin, nicht dem ersten Blicke durchsichtig zu sein. Wenn der Geist reden will, muss das Glück ins Nebenzimmer gehn oder mindestens sich die Ohren zuhalten; sonst ertappt es uns auf einem Undank.
\par

Die Laster verdienen Tadel, weil sie den Menschen nicht zur Meisterschaft im Bösen reifen lassen.
\par

Schlimm, wenn wir durch Alles hindurch müssten, um darüber hinaus zu sein!
\par

Liebe wäre der Schwarzrock, Zeugung zu exorciren!
\par

Des Nöthigsten kann man sich nicht versichern, wie die Lunge nicht erst den Sauerstoffgehalt der Luft bestimmen lassen kann, ehe sie sich zum Athmen entschliesst: sie muss athmen und entweder davonkommen oder ersticken. Sein oder Nichtsein ist keine Frage, sondern ein Versuch.
\par

Wer als Autor auf seinem Persönlichen besteht, verdient ein Publicum, das es geradeso hält.
\par

Was jung Blut niederschrieb, dabei denkt sich grauer Kopf etwas.
\par

Leichter ist todtgeschlagen als die Thüre gewiesen.
\par

Die Grösse der Gefahr misst sich an der Absurdität des Ausweges.
\par

Dass Adel verpflichte, soll der Plebejer nicht in den Mund nehmen. Eigenthum ist Diebstahl, so darf der Eigenthümer reden. Wahrheiten (oder Unwahrheiten) dieser Art sind nur erträglich, wenn, wer sie ausspricht, dadurch einbüsst, also auf Seite der Gewährenden, nicht der Fordernden.
\par

Dass du heute Brot issest, danke den Körnern, die dir beim Säen daneben fielen.
\par

In den letzten Folgen kehren Glück und Fehlgriff ihr Vorzeichen um.
\par

Die Asketen vergessen zu warnen, dass man nicht den Schmerz liebgewinne. Von wann an darf ich mir das Leben wieder erleichtern? — Ein Jahr, einen Tag und eine Stunde vor deinem Tode.
\par

Solang das Leben bergauf geht, soll man nicht sprechen; das wissen alle Bergsteiger. Also, einen Grashalm in den Mund, ihr jungen Dichter!
\par

Durch Zwangslagen wird man organisirt, durch Alternativen zurückgebildet; liberale Zeiten vergeuden, was despotische gesammelt.
\par

Die sich Verständigung schulden, verurtheilen einander zu leben; jederzeit reisefertig ist nur der Einzelne.
\par

Schliesslich verlangt noch die Eitelkeit des Verräthers, dass man seinen Verrath schwer verzeihe.
\par

Du gingest mit guten Dingen um; hüte dich, dass sie später nicht bei dir umgehen, als Gespenster deines gelebten Glücks und Mörder deines lebenden.
\par

Sind wir denn Schauspieler? Glauben wir unsere Ächtheit erst den Zuschauern, die an sie glauben? Und ist es nöthig, dass wir selbst an sie glauben? Auch wir sind Zuschauer unser selbst: was liegt daran, dass Zuschauer an uns glauben?
\par

Ärzte sind die Pfaffen des Leibes, wie Priester die Kurpfuscher der Seele: beide Meister darin, zwanzig Gesundheiten von Leib und Seele über einen Leisten zu schlagen und nach einem Schema zu verderben. In der letzten Lebensstunde mag man immerhin sich in die Hand dieser Beiden liefern.
\par

Um ein geartetes Ich wächst die geartete Welt.
\par

Eure trippelnde Schwärmerei nennt sich Erhabenheit? Soll euch die Erde unter den Füssen fort, so tretet nur fester auf!
\par

Metaphysiker sei nicht, aber sei es gewesen.
\par

Animalisches Leben ist Heldenthum an sich. Der späteren Zukunft lebt der höhere Mensch.
\par

Die älteren Jahrgänge Zufall heissen Nothwendigkeit.
\par

Auf des Mannes Tag hat das Weib kein Recht.
\par

Man muss nicht die Oberfläche That, sondern das Centrum Wollen verbieten, wenn es einmal ans Verbieten geht.
\par

Synthetische Urtheile sind die Selbstanalyse des Urtheilenden.
\par

Das vollkommene lyrische Gedicht schildert nicht Stimmung, sondern ist sie; es lässt sich sowenig wie Musik in Sprache übersetzen; es verhält sich zur normalen Sprache wie überhaupt Kunstwirklichkeit zur empirischen Wirklichkeit.
\par

Die höhere Perversität hegt gegen die niedere denselben Abscheu wie diese gegen Natur.
\par

Bei deinen Tyrannen bedanke dich, nicht bei deinen Freiheiten.
\par

Der Gott wird Bestie, dem Opfer rauchen.
\par

„Wir sind“ ist die Meinung aller Wir, und die Zeit giebt allen Recht.
\par

Der beschämendste Reisebegleiter der That ist — der Thäter.
\par

Wer mich treffen kann, schuldet mir einen Pfeil.
\par

Lass den Gaul nicht wissen, dass du auch Weizen säest.
\par

Von Künstlern und Autoren ist es eigentlich zuviel verlangt, dass sie auch persönlich noch gute Gesellschaft sein sollen.
\par

Unhöflichkeit des Autors, der dem Leser auch nicht Ein Misslungenes vorzusetzen wagt! Man soll einiges dem guten Willen überlassen und den Sieg nicht auf der ganzen Linie erzwingen wollen. Portraittreue Schilderung verräth Mangel an Associationen; wer viel erfahren hat, wird seine Erfahrungen mischen.
\par

Formvollendung, nichts Geringeres, soll der Dichter von sich verlangen, um beim späteren Rückblick auf das Geschaffene über eine bittere Empfindung Herr zu werden: über die Scham nach einem Geständniss.
\par

